% =============================================================================
%  Econometría del modelo de velocidad de arriendo
%  Assetplan · Proyecto "Predicción de tiempo de arriendo"
%
%  Compilar con pdfLaTeX (dos pasadas para las referencias cruzadas).
%  En Overleaf: Menu -> Compiler -> pdfLaTeX.  No requiere bibtex.
% =============================================================================
\documentclass[11pt,a4paper]{article}

\usepackage[utf8]{inputenc}
\usepackage[T1]{fontenc}
\usepackage[spanish,es-tabla,es-noquoting]{babel}
\usepackage{lmodern}
\usepackage{microtype}

\usepackage{amsmath,amssymb,amsthm}
\usepackage{booktabs}
\usepackage{array}
\usepackage{geometry}
\usepackage{enumitem}
\usepackage{xcolor}
\usepackage[hidelinks]{hyperref}

\geometry{margin=2.8cm}

% --- Colores sobrios para marcar la dirección de cada sesgo ------------------
\definecolor{sesgoAbajo}{HTML}{A8452C}
\definecolor{sesgoArriba}{HTML}{8A6410}
\definecolor{neutro}{HTML}{1C6F55}
\newcommand{\dn}{\textcolor{sesgoAbajo}{$\downarrow$}}
\newcommand{\up}{\textcolor{sesgoArriba}{$\uparrow$}}
\newcommand{\nt}{\textcolor{neutro}{$\circ$}}

% --- Entornos ---------------------------------------------------------------
\theoremstyle{definition}
\newtheorem{resultado}{Resultado}
\newtheorem{observacion}{Observación}

% --- Atajos de notación -----------------------------------------------------
\newcommand{\E}{\mathbb{E}}
\newcommand{\Prob}{\mathbb{P}}
\newcommand{\Var}{\operatorname{Var}}
\newcommand{\Cov}{\operatorname{Cov}}
\newcommand{\plim}{\operatorname*{plim}}
\newcommand{\dd}{\,\mathrm{d}}

\title{\bfseries Econometría del modelo de velocidad de arriendo\\[2pt]
\large Especificación, derivadas e inventario de sesgos}
\author{Proyecto \emph{Predicción de tiempo de arriendo}\\ Assetplan}
\date{Septiembre de 2026}

\begin{document}
\maketitle

\begin{abstract}
\noindent
Nota metodológica del modelo de duración de vacancia de unidades de arriendo.
Se presenta la especificación \emph{accelerated failure time} log-normal, se
derivan explícitamente los efectos marginales sobre días y sobre la probabilidad
publicada, y se documenta el inventario completo de sesgos: siete corregidos
dentro del modelo y tres que quedan vivos. Los dos sesgos vivos que afectan al
coeficiente de precio atenúan hacia cero, de modo que la estimación es
plausiblemente una cota inferior del efecto del sobreprecio. Muestra de
estimación: $n=3.889$ unidades en 14 cohortes mensuales, con
$\hat{\sigma}=1{,}335$.
\end{abstract}

\tableofcontents
\bigskip

% =============================================================================
\section{Notación}
% =============================================================================

\begin{table}[h]
\centering
\begin{tabular}{@{}ll@{}}
\toprule
Símbolo & Significado \\
\midrule
$T$ & días desde LPA (\emph{lista para arrendar}) hasta el cierre comercial \\
$t_i$ & tiempo observado, $\min(T_i,\ c_i)$ con $c_i$ el seguimiento disponible \\
$d_i$ & indicador de evento: 1 si se observó el arriendo, 0 si está censurada \\
$x_i$ & vector de atributos conocidos el día de la publicación \\
$g_i$ & \emph{gap} de precio: $\log(\text{precio}) - \log(\text{precio esperado})$ \\
$S(t)$ & función de supervivencia, $\Prob(T>t)$ \\
$h(t)$ & tasa de riesgo (\emph{hazard}) \\
$\phi,\ \Phi$ & densidad y acumulada de la normal estándar \\
$\lambda(z)$ & razón de Mills inversa, $\phi(z)/\bigl(1-\Phi(z)\bigr)$ \\
$\sigma$ & desviación estándar del logaritmo de la duración \\
\bottomrule
\end{tabular}
\end{table}

% =============================================================================
\section{El objeto de estimación: la tasa, no el promedio}
% =============================================================================

El objeto natural no es $\E[T]$ sino la \textbf{tasa de riesgo}: la velocidad
instantánea a la que una unidad vacante se coloca, condicional a seguir vacante,

\begin{equation}
h(t\mid x)=\lim_{\Delta\to 0}
\frac{\Prob\bigl(t\le T<t+\Delta \mid T\ge t,\ x\bigr)}{\Delta}
=\frac{f(t\mid x)}{S(t\mid x)}
=-\frac{\dd}{\dd t}\log S(t\mid x).
\label{eq:hazard}
\end{equation}

Hay dos razones, ambas estructurales. La primera es que el $20{,}0\%$ de la
muestra no tiene desenlace: al cierre de los datos seguía vacante. Un promedio
calculado sobre las unidades que sí cerraron no estima $\E[T]$ sino
$\E[T\mid T\le c]$, que es una cantidad distinta. La segunda es que la tasa es
lo que la operación decide: cuánta atención asignar hoy a una unidad depende de
la probabilidad de que se coloque en los próximos días, no del promedio
histórico.

Integrando \eqref{eq:hazard} se recupera todo lo demás,
$S(t)=\exp\bigl(-\int_0^t h(u)\dd u\bigr)$, y la probabilidad de colocarse antes
de un horizonte $h$ es $1-S(h)$.

% =============================================================================
\section{Especificación: AFT log-normal}
% =============================================================================

La especificación es una regresión lineal sobre el logaritmo de la duración
---\emph{accelerated failure time}, porque los regresores aceleran o frenan el
reloj del proceso---,

\begin{equation}
\log T_i = x_i'\beta + \sigma\varepsilon_i,
\qquad \varepsilon_i \sim N(0,1),
\label{eq:aft}
\end{equation}

de donde $\log T\mid x \sim N(x'\beta,\ \sigma^2)$. Las tres cantidades que el
modelo entrega salen de una única estimación:

\begin{equation}
S(t\mid x)=1-\Phi(z),
\qquad
\Prob(T\le h\mid x)=\Phi\!\left(\frac{\log h-x'\beta}{\sigma}\right),
\qquad
z\equiv\frac{\log t-x'\beta}{\sigma}.
\label{eq:curva}
\end{equation}

Esa es la ventaja práctica sobre un modelo binario: una regresión logística
ajustada sobre ``¿se arrendó en 20 días?'' no contiene información sobre los 30
días. El AFT devuelve la curva completa para cualquier horizonte sin reajustar.

\begin{observacion}[Por qué log-normal y no Weibull]
La distribución Weibull impone una tasa de riesgo monótona: siempre creciente o
siempre decreciente. La evidencia empírica lo contradice: Gdakowicz, Putek-Szeląg
y Bas \cite{gdakowicz2023} documentan una caída pronunciada de la tasa después de
7--8 semanas en el mercado polaco de arriendo. La log-normal admite una tasa no
monótona, que es la forma esperada cuando la demanda inicial se agota. Un modelo
de riesgo en tiempo discreto con la forma base completamente libre ---\emph{spline}
lineal por tramos en $\log(\text{semana})$, siguiendo a Deng, Gabriel y Nothaft
\cite{deng2003}--- alcanza esencialmente el mismo error integrado: 0,195 contra
0,196.
\end{observacion}

% =============================================================================
\section{Efectos marginales}
% =============================================================================

Hay tres derivadas y responden preguntas distintas.

\subsection{Sobre los días}

Directamente de \eqref{eq:aft}, puesto que $\E[\log T\mid x]=x'\beta$:

\begin{equation}
\frac{\partial\,\E[\log T\mid x]}{\partial x_k}=\beta_k,
\qquad
T_{\mathrm{med}}(x)=e^{x'\beta}
\ \Longrightarrow\
\frac{\partial T_{\mathrm{med}}}{\partial x_k}=\beta_k\,T_{\mathrm{med}}.
\label{eq:dias}
\end{equation}

Un coeficiente es una \textbf{semi-elasticidad}: $\beta_k=-0{,}10$ corresponde a
aproximadamente un $10\%$ menos de días. El efecto \emph{en días} no es
constante, sino proporcional al nivel: los mismos $-0{,}10$ equivalen a 3 días en
una unidad que se coloca en 30 y a 9 días en una que tarda 90. Comunicar
``ahorro de días'' sin declarar la base es engañoso.

Cuando el regresor ya está en logaritmos ---como el precio--- el coeficiente es
una \textbf{elasticidad}: variación porcentual de los días ante una variación
porcentual del precio.

\subsection{Sobre la probabilidad publicada}

El entregable no son días sino $\Prob(T\le h)$. Derivando \eqref{eq:curva} por la
regla de la cadena, con $\partial z/\partial x_k=-\beta_k/\sigma$:

\begin{equation}
\boxed{\ \frac{\partial\,\Prob(T\le h\mid x)}{\partial x_k}
=\phi(z)\,\frac{\partial z}{\partial x_k}
=-\,\frac{\phi(z)\,\beta_k}{\sigma}\ }
\label{eq:dprob}
\end{equation}

\begin{resultado}[Dónde tiene efecto la palanca de precio]
El efecto marginal no es constante: lo modula $\phi(z)$, que alcanza su máximo en
$z=0$, es decir cuando la mediana predicha de la unidad coincide exactamente con
el horizonte de reporte. Una rebaja de precio mueve mucho la probabilidad de
colocar en 20 días para una unidad que el modelo estima en torno a 20 días, y
casi nada para una unidad que el modelo estima en 5 o en 90: en los extremos
$\phi(z)$ se aplasta.

La implicación operativa es contraintuitiva: \textbf{la intervención de precio
rinde en el centro de la distribución, no en la cola}. El decil más lento está
tan lejos del horizonte que el precio no lo rescata; allí el problema es de otra
naturaleza.
\end{resultado}

\subsection{Sobre la dispersión}

\begin{equation}
\frac{\partial\,\Prob(T\le h\mid x)}{\partial \sigma}
=\phi(z)\,\frac{z}{\sigma}.
\label{eq:dsigma}
\end{equation}

Un $\sigma$ mayor empuja todas las probabilidades hacia $1/2$ y aplana las
diferencias entre unidades. Con $\hat{\sigma}=1{,}335$ el intervalo de una
desviación estándar corresponde a un factor $e^{1{,}335}\approx 3{,}8$ en días.
Ese número explica por completo que el error absoluto mediano en días sea de
aproximadamente 20 contra una mediana de 29, y es la razón por la cual el
entregable es una probabilidad y un decil, nunca un día puntual.

% =============================================================================
\section{Verosimilitud censurada y ecuaciones de score}
% =============================================================================

Una unidad arrendada aporta su densidad; una unidad todavía vacante aporta la
supervivencia, es decir la información ``tardó \emph{al menos} esto'':

\begin{equation}
\ell(\beta,\sigma)=\sum_i
d_i\bigl[\log\phi(z_i)-\log\sigma\bigr]
+(1-d_i)\log\bigl(1-\Phi(z_i)\bigr),
\qquad z_i=\frac{\log t_i-x_i'\beta}{\sigma}.
\label{eq:loglik}
\end{equation}

Derivando, y usando $\partial z_i/\partial\beta = -x_i/\sigma$,
$\partial z_i/\partial\sigma = -z_i/\sigma$, junto con
$\frac{\dd}{\dd z}\log\phi(z)=-z$ y
$\frac{\dd}{\dd z}\log\bigl(1-\Phi(z)\bigr)=-\lambda(z)$:

\begin{align}
\frac{\partial \ell}{\partial \beta}
&=\frac{1}{\sigma}\sum_i
\Bigl[\,d_i\,z_i+(1-d_i)\,\lambda(z_i)\,\Bigr]x_i,
\label{eq:scorebeta}\\[4pt]
\frac{\partial \ell}{\partial \sigma}
&=\frac{1}{\sigma}\sum_i
\Bigl[\,d_i\,(z_i^{2}-1)+(1-d_i)\,z_i\,\lambda(z_i)\,\Bigr].
\label{eq:scoresigma}
\end{align}

\begin{resultado}[Qué aporta cada observación]
Una observación con evento contribuye su residuo estandarizado $z_i$, igual que
en mínimos cuadrados. Una observación censurada contribuye $\lambda(z_i)$, que es
el \textbf{residuo esperado condicional a superar el umbral observado}. La
censura no se descarta: se imputa con la esperanza condicional que impone el
propio modelo. Esa es toda la diferencia respecto de descartar las 780
observaciones censuradas.
\end{resultado}

\begin{observacion}[Penalización \emph{ridge}]
Lo que se maximiza en la práctica es $\ell(\beta,\sigma)-\lambda\|\beta_{-0}\|^2$,
de modo que \eqref{eq:scorebeta} incorpora un término adicional $-2\lambda\beta$.
Eso introduce sesgo hacia cero de manera deliberada, a cambio de varianza. Con 43
parámetros y aproximadamente 945 eventos el número de eventos por parámetro es
22, holgado según los criterios habituales; el \emph{shrinkage} responde a la
dispersión del diseño ---edificios con pocas unidades--- más que a escasez de
muestra. Sigue la recomendación de van Smeden \emph{et al.} \cite{vansmeden2019}.
\end{observacion}

% =============================================================================
\section{Sesgo I: censura}
% =============================================================================

Es el sesgo de mayor magnitud y el modelo lo corrige. Supóngase el procedimiento
intuitivo: regresión por mínimos cuadrados de $\log T$ sobre las unidades ya
arrendadas. Eso estima $\E[\log T\mid x,\ T\le c]$. Para una normal truncada por
arriba, con $a\equiv(\log c-x'\beta)/\sigma$ y usando
$\E[Z\mid Z\le a]=-\phi(a)/\Phi(a)$ para $Z\sim N(0,1)$:

\begin{equation}
\E\bigl[\log T \mid x,\ T\le c\bigr]
= x'\beta-\sigma\,\frac{\phi(a)}{\Phi(a)}.
\label{eq:truncada}
\end{equation}

El término omitido, $-\sigma\phi(a)/\Phi(a)$, es siempre negativo: el estimador
ingenuo \textbf{siempre predice más rápido de lo real}. Pero lo grave no es el
signo sino que $a$ depende de $x'\beta$. El término omitido está correlacionado
con los regresores, de modo que no es un desplazamiento del intercepto que pueda
ignorarse: \textbf{sesga todo el vector $\beta$}.

El problema se agrava justamente donde más importa. Cuanto más corto el
seguimiento ---las cohortes recientes--- menor es $a$, mayor $\phi(a)/\Phi(a)$ y
más fuerte el sesgo. En la cohorte 2026-08 la censura alcanzó el $73{,}9\%$.

\begin{resultado}[Magnitud medida]
Mediana de días predicha por el AFT censurado: $47{,}6$. Por mínimos cuadrados
descartando las censuradas: $23{,}2$. El sesgo es de $-24{,}4$ días, más de la
mitad del valor estimado.
\end{resultado}

% =============================================================================
\section{Sesgo II: riesgo competitivo}
% =============================================================================

El $18{,}2\%$ de las unidades no se arrienda: se da de baja. Eso \textbf{no es
censura}: no es ``todavía no se sabe'', es un desenlace observado y definitivo.
Existen dos causas con tasas específicas $h_1$ (arriendo) y $h_2$ (baja), y la
supervivencia depende de ambas:

\begin{equation}
S(t)=\exp\!\left(-\int_0^t\bigl[h_1(u)+h_2(u)\bigr]\dd u\right),
\qquad
F_1(t)=\int_0^t S(u^-)\,h_1(u)\dd u.
\label{eq:cif}
\end{equation}

$F_1$ es la \textbf{incidencia acumulada}: la probabilidad de que la unidad se
arriende antes de $t$ en el mundo real, donde puede darse de baja antes. Si en
cambio se trata la baja como censura y se calcula
$1-\exp(-\int_0^t h_1)$, se está estimando la probabilidad de arriendo en un
mundo contrafactual sin bajas. La desigualdad tiene una sola dirección:

\begin{equation}
1-\exp\!\left(-\int_0^t h_1(u)\dd u\right)
\ \ge\
\int_0^t S(u^-)h_1(u)\dd u
\ =\ F_1(t),
\label{eq:desigualdad}
\end{equation}

cuya demostración es inmediata: $\exp\bigl(-\int(h_1+h_2)\bigr)\le
\exp\bigl(-\int h_1\bigr)$ porque $h_2\ge 0$. Tratar la baja como censura
\textbf{siempre sobreestima} la probabilidad de arriendo. La formulación de
referencia para regresión sobre la incidencia acumulada es Fine y Gray
\cite{finegray1999}.

\begin{observacion}[Dónde muerde]
A 20 días la masa de bajas es pequeña: $\Prob(\text{arrendada})=0{,}263$ frente a
$\Prob(\text{baja})=0{,}024$. Pero la baja no está repartida uniformemente: se
concentra en los segmentos lentos, con un $43\%$ de censura en comunas como
Cerrillos y San Joaquín. El sesgo es mayor exactamente donde el modelo debe
advertir un problema.
\end{observacion}

% =============================================================================
\section{El precio de referencia y sus dos sesgos}
% =============================================================================

El precio en nivel no sirve: entre deciles extremos, las unidades caras se
colocan $10{,}9$ días \emph{más rápido}, porque las comunas caras son las
líquidas. El precio solo informa relativo a lo que la unidad debería costar. De
ahí la estructura en dos etapas, que sigue a Andersson \cite{andersson2025}:

\begin{align}
\text{(i)}\quad & \log P_i = w_i'\gamma + u_i, \label{eq:etapa1}\\
\text{(ii)}\quad & \log T_i = \alpha\,\hat{u}_i + x_i'\beta + \sigma\varepsilon_i.
\label{eq:etapa2}
\end{align}

La primera etapa es un modelo hedónico del precio; el residuo $\hat{u}$ es el
\emph{gap}. Por construcción $\hat{u}$ es ortogonal a $w$ en la muestra de
ajuste: es \textbf{la componente del precio que los atributos no explican}, es
decir la parte discrecional de la decisión de \emph{pricing}. La justificación
teórica del canal es la búsqueda competitiva de Díaz y Jerez \cite{diazjerez2010}:
publicando precios más bajos, el vendedor atrae más visitas y coloca más rápido.

\subsection{Regresor generado}

$\hat{u}$ no se observa, se estima. Eso rompe dos cosas distintas, y conviene no
confundirlas porque tienen remedios distintos.

\paragraph{(a) Inferencia.} Los errores estándar de la segunda etapa están
subestimados porque ignoran la incertidumbre de $\hat{\gamma}$: es el problema de
Pagan \cite{pagan1984}. La corrección es aplicar \emph{bootstrap} al procedimiento
completo ---reestimando \eqref{eq:etapa1} dentro de cada réplica--- o la corrección
analítica de Murphy y Topel \cite{murphytopel1985}. Es lo que hacen Allen,
Rutherford y Thomson \cite{allen2009} sobre 20.131 arriendos de Dallas con esta
misma especificación en dos etapas.

\paragraph{(b) Sesgo.} Si $\hat{u}=u-\xi$ con $\xi=\log\hat{P}-\log P^{*}$ el error
de estimación de la primera etapa, se trata de error de medición clásico en el
regresor y

\begin{equation}
\plim\ \hat{\alpha}
=\alpha\cdot\underbrace{\frac{\Var(u)}{\Var(u)+\Var(\xi)}}_{\textstyle\lambda}\ <\ \alpha .
\label{eq:atenuacion}
\end{equation}

Atenuación hacia cero: cuanto peor el modelo de precio, más se diluye el efecto
estimado del \emph{gap}. $\lambda$ es la \emph{razón de confiabilidad}
\cite{fuller1987}.

\begin{observacion}[El \emph{bootstrap} no corrige (b)]
Un \emph{bootstrap} reproduce en cada réplica el sesgo que tenga el estimador. Si
$\hat{\alpha}$ está atenuado, las $B$ réplicas están atenuadas y el resultado es un
intervalo angosto alrededor del número equivocado: precisión sobre algo sesgado. El
resultado de Pagan es sobre la \textbf{varianza}, no sobre el punto. Allen et al.
tampoco corrigen la atenuación.
\end{observacion}

\subsection{Identificación de $\Var(\xi)$ y corrección de la atenuación}

Para aplicar \eqref{eq:atenuacion} hace falta $\Var(\xi)$, y el obstáculo es que
\textbf{no se identifica con el residuo del hedónico}: ese residuo \emph{es}
$\hat{u}$, y bajo $\xi\perp u$ ya trae las dos varianzas sumadas,

\begin{equation}
\Var(\hat{u})=\Var(u)+\Var(\xi).
\label{eq:varsuma}
\end{equation}

Hace falta una fuente de variación que mueva $\hat{P}$ dejando $P^{*}$ fijo. El
propio \emph{bootstrap} de (a) la provee. Para una unidad \textbf{fija} $i$, se
recalcula el \emph{gap} en cada réplica con el hedónico reajustado sobre la
remuestra. Como $u_i$ no cambia entre réplicas,

\begin{equation}
\hat{u}_i^{(b)}-\hat{u}_i
=-\bigl(\log\hat{P}^{(b)}(w_i)-\log\hat{P}(w_i)\bigr),
\label{eq:gapreplica}
\end{equation}

de modo que la dispersión entre réplicas aísla exactamente la variabilidad de la
primera etapa y

\begin{equation}
\widehat{\Var}(\xi)=\frac{1}{n}\sum_{i=1}^{n}\Var_b\!\left(\hat{u}_i^{(b)}\right),
\qquad
\hat{\lambda}=1-\frac{\widehat{\Var}(\xi)}{\Var(\hat{u})},
\qquad
\hat{\alpha}_{\text{corr}}=\frac{\hat{\alpha}}{\hat{\lambda}} .
\label{eq:lambda}
\end{equation}

El mismo factor $1/\hat{\lambda}$ se aplica a los extremos del intervalo. Es una
corrección de primer orden: trata $\hat{\lambda}$ como conocido e ignora su propia
varianza de muestreo, lo que subestima levemente la amplitud del intervalo
corregido.

\begin{resultado}[Qué corrige y qué no]
$\hat{\alpha}/\hat{\lambda}$ sigue siendo una \textbf{cota inferior}, por dos razones
independientes:

\begin{enumerate}[leftmargin=1.6em,itemsep=2pt]
\item El \emph{bootstrap} mide el error de \textbf{estimación} de $\hat{P}$, no el de
\textbf{especificación}. Si el hedónico está mal especificado, $\hat{P}$ es sesgado
para $P^{*}$ y esa componente de $\xi$ no aparece en la dispersión entre réplicas
---es común a todas ellas---. Luego $\widehat{\Var}(\xi)$ es una subestimación,
$\hat{\lambda}$ una sobrestimación, y la corrección se queda corta.
\item Corrige la atenuación por medición, no la endogeneidad de la sección
siguiente, que empuja en la misma dirección.
\end{enumerate}
\end{resultado}

Una alternativa sin este límite es SIMEX \cite{cookstefanski1994}: se inyecta ruido
adicional de varianza conocida $\theta\,\widehat{\Var}(\xi)$ para $\theta=0{,}5,1,1{,}5,2$,
se observa cómo decae $\hat{\alpha}(\theta)$ y se extrapola a $\theta=-1$. Es más
robusto a la forma funcional del decaimiento, y también más caro: multiplica por
cinco el número de ajustes.

\subsection{Endogeneidad: el sesgo que no se corrige}

La ecuación estructural es $\log T=\alpha g+x'\beta+\nu$, donde $\nu$ contiene
todo lo que hace deseable a una unidad y no se observa: orientación, vista, luz,
estado de las terminaciones. Pero \textbf{quien fija el precio sí observa parte
de $\nu$}. La regla de \emph{pricing} es $g=f(x,\nu,\eta)$, y en consecuencia

\begin{equation}
\plim\ \hat{\alpha}=\alpha+\frac{\Cov(g,\nu)}{\Var(g\mid x)} .
\label{eq:endogeneidad}
\end{equation}

El signo del segundo término admite un argumento. Dos historias plausibles, y
ambas apuntan en la misma dirección:

\begin{itemize}[leftmargin=1.4em,itemsep=2pt]
\item La unidad es mejor de lo que registran los datos ($\nu$ bajo, se coloca
rápido) y quien fija el precio lo advierte y cobra más ($g$ alto)
$\Rightarrow \Cov(g,\nu)<0$.
\item La unidad es peor de lo que registran los datos ($\nu$ alto, lenta) y se le
asigna un precio de castigo ($g$ bajo) $\Rightarrow \Cov(g,\nu)<0$ también.
\end{itemize}

\begin{resultado}[Dirección del sesgo de endogeneidad]
Con $\Cov(g,\nu)<0$ y $\alpha>0$, la estimación \textbf{subestima} el efecto del
sobreprecio. Sumado a la atenuación de \eqref{eq:atenuacion}, que opera en la
misma dirección, $\hat{\alpha}$ es plausiblemente una \textbf{cota inferior} del
daño real de publicar caro.

Es un argumento de signo, no una demostración: descansa en que la regla de
\emph{pricing} no asigne precios altos a unidades sistemáticamente peores. Pero
es la dirección conservadora, lo que significa que el modelo no exagera el efecto
del precio.
\end{resultado}

Corregirlo exigiría variación exógena de precios: un experimento, o un cambio de
política que haya movido precios por razones ajenas a la unidad. Es el problema
que formaliza Stein \cite{stein1995}: el precio de lista es una elección
estratégica del vendedor, no un tratamiento asignado. \textbf{Reducir el precio de
una unidad no produce mecánicamente el cambio de velocidad que sugiere el
coeficiente.}

\begin{observacion}[Causalidad inversa dentro de la ventana]
El $37{,}3\%$ de las unidades tuvo una rebaja durante la vacancia, y esas tardan
53 días de mediana frente a 24. Eso \emph{no} es evidencia de que rebajar no
sirva: se rebaja \emph{porque} no se arrienda. El \emph{gap} se mide con el precio
del día 0 precisamente para evitar esta contaminación, pero la protección es
parcial: si el precio inicial ya anticipaba un problema, la contaminación entra
igualmente.
\end{observacion}

% =============================================================================
\section{Sesgo IV: codificación del efecto de edificio}
% =============================================================================

Hay 796 edificios y 3.889 unidades: como variables indicadoras es inviable. Se
utiliza la media del objetivo por edificio, encogida hacia la media global ---un
estimador de Bayes empírico---:

\begin{equation}
\mathrm{TE}_j=\frac{n_j\,\bar{y}_j+m\,\bar{y}}{n_j+m}
\qquad\Longrightarrow\qquad
\E[\mathrm{TE}_j]-\theta_j=\frac{m}{n_j+m}\,(\bar{y}-\theta_j).
\label{eq:te}
\end{equation}

Sesgo deliberado hacia la media global, que se desvanece a medida que $n_j$
crece, a cambio de reducir la varianza por un factor $\bigl(n_j/(n_j+m)\bigr)^2$.
El $m$ óptimo en error cuadrático medio es aproximadamente el cociente entre la
varianza intra-edificio y la varianza entre edificios; en la implementación está
fijado en 20.

\begin{observacion}[El sesgo peligroso es el otro]
Si el código del edificio $j$ se calcula empleando el $y_i$ de la propia
observación $i$, el regresor queda correlacionado con el error por construcción y
el ajuste dentro de muestra se infla de manera espectacular. Es la forma más
común de inflar artificialmente un AUC. La codificación se calcula fuera de
pliegue ---cinco particiones internas dentro del conjunto de entrenamiento---
precisamente por esto. Nótese además que para un edificio nuevo sin historial el
código colapsa al valor previo global: el modelo pierde su variable más fuerte
exactamente donde más se la querría.
\end{observacion}

% =============================================================================
\section{Sesgo V: el esquema de validación}
% =============================================================================

La validación cruzada aleatoria estima el error esperado bajo el supuesto de que
las observaciones son extracciones independientes de una distribución fija. Aquí
ese supuesto falla por dos vías simultáneas.

\textbf{Correlación intra-edificio.} Dos unidades de la misma torre comparten
ubicación, administración y política de precio. El tamaño de muestra efectivo se
aproxima al número de edificios más que al de unidades, y con partición aleatoria
el mismo edificio aparece en entrenamiento y en prueba.

\textbf{No estacionariedad.} La distribución conjunta se desplaza entre cohortes:
las visitas de la primera semana oscilan entre 4,7 y 13,3 según el mes.

\begin{table}[h]
\centering
\begin{tabular}{@{}lcl@{}}
\toprule
Esquema & AUC@20 & Supuesto que emplea \\
\midrule
Validación cruzada aleatoria & 0,767 & observaciones i.i.d.: falso por ambas vías \\
\emph{GroupKFold} por edificio & 0,731 & corrige el conglomerado, no la deriva \\
\textbf{\emph{Forward} por cohorte} & \textbf{0,705} & reproduce el uso real: pasado $\to$ mes siguiente \\
\bottomrule
\end{tabular}
\caption{Efecto del esquema de validación sobre el desempeño estimado.}
\end{table}

La descomposición es limpia: $0{,}036$ de fuga por conglomerado y $0{,}026$ de
deriva temporal. Los $0{,}062$ conjuntos son el margen por el cual un proyecto
validado con partición aleatoria promete un desempeño que no obtendrá.

% =============================================================================
\section{Calibración}
% =============================================================================

La regresión de calibración de Cox contrasta el \emph{nivel} de la probabilidad,
no su orden:

\begin{equation}
\operatorname{logit}\Prob(Y=1\mid \hat{p})=a+b\cdot\operatorname{logit}(\hat{p}).
\label{eq:calib}
\end{equation}

Con $b=1$ y $a=0$ la probabilidad significa lo que declara. Un $b<1$ indica
predicciones sobredispersas: los extremos exageran. Las pendientes medidas fueron
$0{,}487$ para la logística, $0{,}717$ para el AFT y $0{,}725$ para el modelo de
riesgos competitivos; la primera lo suficientemente baja como para que sus
probabilidades no sean publicables.

El escalamiento de Platt \cite{platt1999} lo corrige reajustando $a$ y $b$, y esa
corrección \textbf{también se estima hacia adelante}: sobre las cohortes
anteriores, aplicada a la siguiente. Ajustarla sobre las mismas predicciones que
después se evalúan produce una pendiente de $1{,}000$ por construcción y carece
de contenido.

\begin{observacion}[Dos verificaciones]
La transformación es monótona, de modo que el AUC \emph{dentro} de una cohorte no
puede cambiar; ésa es la comprobación de que la implementación es correcta. Y si
la pendiente estimada resulta negativa, el mapeo invierte el orden y el
\emph{score} recalibrado queda peor que el original: en ese caso la cohorte se
deja sin recalibrar.
\end{observacion}

% =============================================================================
\section{El techo estocástico}
% =============================================================================

Parte del error no es eliminable con ninguna variable adicional, porque el
desenlace \emph{es} un tiempo de espera aleatorio. Para el error cuadrático de la
probabilidad la descomposición es exacta, con $\pi$ la probabilidad verdadera:

\begin{equation}
\E\bigl[(Y-\hat{p})^2\bigr]
=\underbrace{\E\bigl[\pi(1-\pi)\bigr]}_{\text{irreducible}}
+\underbrace{\E\bigl[(\pi-\hat{p})^2\bigr]}_{\text{error del modelo}} .
\label{eq:brier}
\end{equation}

Para el ordenamiento el resultado es igual de nítido. Consideremos dos unidades
con medias $\mu_1$ y $\mu_2$, cada una con su propio choque log-normal:

\begin{equation}
\Prob(T_1<T_2\mid \mu_1,\mu_2)
=\Prob\bigl(\sigma(\varepsilon_1-\varepsilon_2)<\mu_2-\mu_1\bigr)
=\Phi\!\left(\frac{\mu_2-\mu_1}{\sigma\sqrt{2}}\right).
\label{eq:concordancia}
\end{equation}

\begin{resultado}[Lo que implica \eqref{eq:concordancia}]
La concordancia máxima alcanzable depende \textbf{únicamente del cociente entre la
dispersión de las medias predichas y $\sigma$}. No del algoritmo empleado ni del
tamaño de muestra. Si $\sigma$ es grande respecto de cuánto se separan las
unidades entre sí, ni un modelo perfecto ordena bien.
\end{resultado}

\begin{table}[h]
\centering
\begin{tabular}{@{}lccc@{}}
\toprule
Métrica & Observado & Techo & \% capturado \\
\midrule
Índice de concordancia & 0,685 & 0,713 & 87\% \\
AUC@20 & 0,725 & 0,763 & 86\% \\
\bottomrule
\end{tabular}
\caption{Desempeño observado frente al máximo alcanzable, estimado por simulación
desde el propio modelo ajustado.}
\end{table}

La simulación de sensibilidad lo confirma: con $\sigma$ reducido a la mitad el
techo del índice de concordancia subiría a $0{,}811$. El límite lo impone el ruido
del proceso, no el algoritmo. Andersson \cite{andersson2025} realiza el mismo
cálculo con 20.000 unidades en Japón y obtiene un techo de aproximadamente 0,66
frente a un índice logrado de 0,63. El índice utilizado es el de Harrell
\emph{et al.} \cite{harrell1982}, en la variante consistente bajo censura de Uno
\emph{et al.} \cite{uno2011}.

\begin{observacion}[Qué no dice este techo]
Se calcula bajo el supuesto de que el modelo está correctamente especificado, es
decir que las $\hat{p}$ son las verdaderas. Si faltan variables, las
probabilidades verdaderas se separan más entre unidades y \textbf{el techo real es
más alto}. La conclusión correcta no es ``no se puede predecir mejor'' sino ``no
se puede predecir mejor \emph{con esta información}''.
\end{observacion}

% =============================================================================
\section{Inventario de sesgos}
% =============================================================================

\begin{table}[h]
\centering
\small
\begin{tabular}{@{}p{4.6cm}cp{2.4cm}p{5.4cm}@{}}
\toprule
Sesgo & Dir. & Magnitud & Estado \\
\midrule
Censura & \dn & $-24{,}4$ días & corregido: verosimilitud \eqref{eq:loglik} \\
Riesgo competitivo & \up & 0,024 a 20 d & corregido: incidencia acumulada \eqref{eq:cif} \\
Fuga en la codificación & \up & --- & corregido: codificación fuera de pliegue \\
Validación aleatoria & \up & $+0{,}062$ & corregido: \emph{forward} por cohorte \\
Sesgo de longitud & \up & 73,9\% censura & corregido: se excluyen cohortes con observabilidad $<0{,}90$ \\
Descalibración & \up & $b=0{,}717$ & corregido: Platt hacia adelante \eqref{eq:calib} \\
\emph{Shrinkage ridge} & \dn & EPV $=22$ & deliberado: se acepta por varianza \\
\addlinespace
Reg.\ generado: inferencia & --- & $\mathrm{se}_{\text{comp}}/\mathrm{se}_{\text{ing}}$ & corregido: \emph{bootstrap} de ambas etapas \\
Reg.\ generado: atenuación & \dn & $\hat{\lambda}$, ec.\ \eqref{eq:lambda} & \textbf{parcial}: $\hat{\alpha}/\hat{\lambda}$ sigue siendo cota inferior \\
Endogeneidad del precio & \dn & --- & \textbf{vivo}: requiere variación exógena \\
Selección de muestra & \nt & 91\% de las LPA & \textbf{vivo}: se declara, no se corrige \\
\bottomrule
\end{tabular}
\caption{Inventario de sesgos. \dn\ atenúa o subestima; \up\ infla o sobreestima;
\nt\ afecta validez externa, no interna.}
\end{table}

Los ocho primeros están incorporados al modelo. La atenuación por regresor
generado está corregida solo en parte: \eqref{eq:lambda} recupera la componente
atribuible al error de estimación de la primera etapa, no la que viene de su
especificación. Los dos últimos no son corregibles con métodos estadísticos sobre
estos datos, y por lo tanto deben declararse en cualquier informe que utilice estos
números.

Los sesgos que siguen afectando al coeficiente de precio ---la atenuación
residual y la endogeneidad--- \textbf{operan ambos hacia cero}. Ese es el hallazgo más tranquilizador del inventario: el modelo no
exagera el efecto del precio, lo subestima. El resultado de la ablación
---eliminar el \emph{gap} cuesta $0{,}052$ de AUC--- es entonces un piso y no un
techo.

El sesgo de selección es de otra naturaleza: no deforma nada dentro de la muestra,
define a quién aplica el modelo. Cubre unidades que alcanzan el estado ``lista
para arrendar'' con historial de precio. Quedan fuera las 1.733 unidades que
alcanzaron el estado ``arrendado'' sin pasar nunca por publicación ---ingresaron
ya ocupadas o fueron asignadas directamente--- y aproximadamente el $9\%$ sin
historial de precio.

% =============================================================================
\section{Síntesis}
% =============================================================================

\begin{enumerate}[leftmargin=1.6em,itemsep=5pt]
\item \textbf{El modelo estima una tasa, no un promedio}, porque una de cada cinco
unidades carece de desenlace y descartarlas sesga la predicción en 24 días.

\item \textbf{Un coeficiente es una semi-elasticidad sobre los días}, y su efecto
sobre la probabilidad publicada es máximo cuando la mediana predicha de la unidad
coincide con el horizonte, según \eqref{eq:dprob}. La palanca de precio rinde en
el centro de la distribución, no en la cola.

\item \textbf{El efecto del precio está subestimado, no exagerado}, porque los
sesgos que permanecen atenúan hacia cero. El \emph{bootstrap} de dos etapas corrige
la inferencia y \eqref{eq:lambda} recupera parte de la atenuación, pero ni uno ni
otro tocan la endogeneidad. Sigue siendo, por tanto, una asociación: para afirmar
que rebajar acelera se requiere variación exógena que estos datos no contienen.
\end{enumerate}

% =============================================================================
\begin{thebibliography}{99}

\bibitem{andersson2025}
Andersson, P. (2025).
\emph{Predicting the vacancy duration of Japanese rental apartments}.
Proceedings of the 39th Annual Conference of the Japanese Society for Artificial
Intelligence, 2O4-OS-21a-03.

\bibitem{deng2003}
Deng, Y., Gabriel, S.\ A.\ y Nothaft, F.\ E. (2003).
Duration of residence in the rental housing market.
\emph{Journal of Real Estate Finance and Economics}, 26(2/3), 267--285.

\bibitem{diazjerez2010}
Díaz, A.\ y Jerez, B. (2010).
\emph{House prices, sales and time on the market: a search-theoretic framework}.
Universidad Carlos III de Madrid, Working Paper Economic Series 10-33.

\bibitem{finegray1999}
Fine, J.\ P.\ y Gray, R.\ J. (1999).
A proportional hazards model for the subdistribution of a competing risk.
\emph{Journal of the American Statistical Association}, 94(446), 496--509.

\bibitem{allen2009}
Allen, M.\ T., Rutherford, R.\ C.\ y Thomson, T.\ A. (2009).
Residential asking rents and time on the market.
\emph{Journal of Real Estate Finance and Economics}, 38(4), 351--365.

\bibitem{cookstefanski1994}
Cook, J.\ R.\ y Stefanski, L.\ A. (1994).
Simulation-extrapolation estimation in parametric measurement error models.
\emph{Journal of the American Statistical Association}, 89(428), 1314--1328.

\bibitem{fuller1987}
Fuller, W.\ A. (1987).
\emph{Measurement Error Models}. Wiley, Nueva York.

\bibitem{gdakowicz2023}
Gdakowicz, A., Putek-Szeląg, E.\ y Bas, M. (2023).
Duration of the rental offer for residential property.
\emph{Real Estate Management and Valuation}, 31(2), 82--91.

\bibitem{gdakowicz2025}
Gdakowicz, A.\ y Putek-Szeląg, E. (2025).
The impact of selected attributes of new residential properties on the timing of
their sale.
\emph{Real Estate Management and Valuation}, 33(3), 60--74.

\bibitem{hanhausman1990}
Han, A.\ y Hausman, J.\ A. (1990).
Flexible parametric estimation of duration and competing risk models.
\emph{Journal of Applied Econometrics}, 5(1), 1--28.

\bibitem{harrell1982}
Harrell, F.\ E., Califf, R.\ M., Pryor, D.\ B., Lee, K.\ L.\ y Rosati, R.\ A. (1982).
Evaluating the yield of medical tests.
\emph{JAMA}, 247(18), 2543--2546.

\bibitem{murphytopel1985}
Murphy, K.\ M.\ y Topel, R.\ H. (1985).
Estimation and inference in two-step econometric models.
\emph{Journal of Business \& Economic Statistics}, 3(4), 370--379.

\bibitem{pagan1984}
Pagan, A. (1984).
Econometric issues in the analysis of regressions with generated regressors.
\emph{International Economic Review}, 25(1), 221--247.

\bibitem{platt1999}
Platt, J. (1999).
Probabilistic outputs for support vector machines and comparisons to regularized
likelihood methods. En \emph{Advances in Large Margin Classifiers}, MIT Press.

\bibitem{stein1995}
Stein, J.\ C. (1995).
Prices and trading volume in the housing market: a model with downpayment effects.
\emph{Quarterly Journal of Economics}, 110(2), 379--406.
(Versión previa: NBER Working Paper 4373, 1993.)

\bibitem{uno2011}
Uno, H., Cai, T., Pencina, M.\ J., D'Agostino, R.\ B.\ y Wei, L.-J. (2011).
On the C-statistics for evaluating overall adequacy of risk prediction procedures
with censored survival data.
\emph{Statistics in Medicine}, 30(10), 1105--1117.

\bibitem{vansmeden2019}
van Smeden, M., Moons, K.\ G.\ M., de Groot, J.\ A.\ H., Collins, G.\ S.,
Altman, D.\ G., Eijkemans, M.\ J.\ C.\ y Reitsma, J.\ B. (2019).
Sample size for binary logistic prediction models: beyond events per variable
criteria.
\emph{Statistical Methods in Medical Research}, 28(8), 2455--2474.

\end{thebibliography}

\end{document}
